% !TeX program = xelatex
\documentclass[UTF8,11pt,fontset=fandol]{ctexart}
\usepackage[a4paper,margin=22mm]{geometry}
\usepackage{amsmath,amssymb,booktabs,longtable}
\usepackage[colorlinks=true,linkcolor=blue]{hyperref}
\setlength{\parindent}{0pt}
\setlength{\parskip}{5pt}
\renewcommand{\arraystretch}{1.3}
\title{乘法与复合函数求导：公式速查}
\author{雪落长安无声}
\date{2026年11月7日增补版}
\begin{document}
\pagestyle{plain}
\maketitle
\section{乘法求导}
各因子须在所求点有相应阶数的导数，但不要求因子本身非零。常数 $c$ 满足 $(cu)'=cu'$，和差满足 $(u\pm v)'=u'\pm v'$。
\begin{align*}
 (uv)'&=u'v+uv',\\
 (uvw)'&=u'vw+uv'w+uvw',\\
 (u_1\cdots u_m)'&=\sum_{k=1}^{m}u_k'\prod_{j\ne k}u_j,\\
 (uv)''&=u''v+2u'v'+uv'',\\
 (uv)'''&=u'''v+3u''v'+3u'v''+uv''',\\
 (uv)^{(n)}&=\sum_{k=0}^{n}\binom nk u^{(k)}v^{(n-k)}.
\end{align*}
这里 $u^{(0)}=u$，二项式系数为 $\binom nk=n!/[k!(n-k)!]$。多因子高阶推广：
\[
 (u_1\cdots u_m)^{(n)}=
 \sum_{\substack{k_1+\cdots+k_m=n\\k_j\ge0}}
 \frac{n!}{k_1!\cdots k_m!}u_1^{(k_1)}\cdots u_m^{(k_m)}.
\]
例如 $(x^2\sin x)'=2x\sin x+x^2\cos x$，不是 $2x\cos x$。
商的求导另为 $(u/v)'=(u'v-uv')/v^2$，要求 $v\ne0$。

\section{复合函数求导}
内函数在所求点可导，外函数在对应内函数值处可导，才能逐层使用链式法则。
\begin{align*}
 [f(u(x))]'&=f'(u(x))u'(x),\\
 [f(g(u(x)))]'&=f'(g(u(x)))g'(u(x))u'(x),\\
 [f(u)]''&=f''(u)(u')^2+f'(u)u'',\\
 [f(u)]'''&=f'''(u)(u')^3+3f''(u)u'u''+f'(u)u'''.
\end{align*}
线性内函数还满足 $[f(ax+b)]^{(n)}=a^nf^{(n)}(ax+b)$（$n\ge1$）。一般内函数不可直接照搬。
\begin{longtable}{@{}p{.23\textwidth}p{.40\textwidth}p{.31\textwidth}@{}}
\toprule 函数&关于 $x$ 的导数&实值区间与条件\\\midrule\endhead
$[u(x)]^\alpha$&$\alpha u^{\alpha-1}u'$&任意实 $\alpha$ 统一取 $u>0$\\
$1/u$&$-u'/u^2$&$u\ne0$\\
$\sqrt{u}$&$u'/(2\sqrt{u})$&此公式要求 $u>0$\\
$e^u$&$e^u u'$&$u$ 可导\\
$a^u$&$a^u\ln a\,u'$&$a>0$；$a=1$ 时导数零\\
$\ln u$&$u'/u$&$u>0$\\
$\ln|u|$&$u'/u$&$u\ne0$\\
$\log_a u$&$u'/(u\ln a)$&$u>0,a>0,a\ne1$\\
$\sin u$&$\cos u\,u'$&弧度制\\
$\cos u$&$-\sin u\,u'$&注意负号\\
$\tan u$&$\sec^2u\,u'$&$\cos u\ne0$\\
$\cot u$&$-\csc^2u\,u'$&$\sin u\ne0$\\
$\sec u$&$\sec u\tan u\,u'$&$\cos u\ne0$\\
$\csc u$&$-\csc u\cot u\,u'$&$\sin u\ne0$\\
$\arcsin u$&$u'/\sqrt{1-u^2}$&$|u|<1$\\
$\arccos u$&$-u'/\sqrt{1-u^2}$&$|u|<1$\\
$\arctan u$&$u'/(1+u^2)$&$u$ 可导\\
$\operatorname{arccot}u$&$-u'/(1+u^2)$&取值为 $(0,\pi)$ 的分支\\
$u^v$&$u^v\bigl(v'\ln u+vu'/u\bigr)$&统一取 $u>0$，$u,v$ 可导\\
\bottomrule
\end{longtable}

表中记 $u=u(x),v=v(x)$；令 $u=x,u'=1$ 就得到基本函数导数。整数幂可按实际定义域放宽 $u>0$ 条件；根式、反三角的特殊端点应直接核查函数，而非把零分母代入公式。

\section{两种规则一起用}
\begin{align*}
 [e^{x^2}\sin(3x)]'&=2xe^{x^2}\sin(3x)+3e^{x^2}\cos(3x),\\
 [\ln(1+x^2)]'&=\frac{2x}{1+x^2},\\
 [\sin(x^2)]''&=2\cos(x^2)-4x^2\sin(x^2).
\end{align*}
检查顺序：先识别乘积或复合结构，再求各层导数，最后检查负号、系数、定义域与分母。
\end{document}
